% GENERATED FILE. Edit canonical lessons, then run npm run generate:books.
% canonical-source-hash: fnv1a64:b3617ecda3704494
% canonical-lessons: ZH-C504-geci, ZH-C504-erge, ZH-W504-ge-song, ZH-C504-guoge, ZH-C504-minge, ZH-C504-yanchang-sing

\chapter{Songs}
\label{ch:zh-songs}
\hlchaptermodality{504}

\begin{chapteropening}
\textbf{By the end of this chapter:} \emph{I can talk about the words of a song, nursery rhymes, the national anthem, folk songs and singing on stage.}

Complete the last lesson of chapter 504: I can talk about the words of a song, nursery rhymes, the national anthem, folk songs and singing on stage.
\end{chapteropening}

\section[gēcí]{gēcí — the words of a song; lyrics}

\label{lesson:ZH-C504-geci}



\begin{quote}
Before the new one: say the Mandarin for to create (art or writing), then the Mandarin for a new year print.
\end{quote}

\subsection*{You'll want to know: gēcí}
\emph{gēcí} — \textquotedblleft{}the words of a song; lyrics\textquotedblright{}.

\begin{quote}
\emph{wǒ tīng le liǎng tiān cái jìzhù gēcí} — It took two days of listening before I knew the words.
\end{quote}

Its characters come later; for now it is a word for the ear.

\begin{cousinweb}
\emph{gēcí} is built from \emph{gē} 'a song' and \emph{cí} 'words', both inherited from Middle Chinese, \emph{ka} and \emph{zi} (Guangyun, in Baxter's transcription), which glosses the first 'the voice of a person singing' and the second 'to speak'.
\end{cousinweb}

\begin{grammarlens}[title={What you've built}]
One more word for saying how it was.
\end{grammarlens}

\subsection*{Your turn}
\begin{itemize}
\raggedright
  \item \emph{Say it:} \emph{gēcí}
  \item \emph{Say it:} \emph{gēcí}, once more
  \item \emph{Say it:} \emph{wǒ tīng le liǎng tiān cái jìzhù gēcí} — It took two days of listening before I knew the words.
  \item \emph{Recall:} say the Mandarin for to create (art or writing), then the Mandarin for a new year print, then say \emph{gēcí} again
\end{itemize}

\begin{tcolorbox}[breakable,colback=teal!4,colframe=teal!35!black,title={Before you move on}]
What does gēcí mean? (\textquotedblleft{}The words of a song; lyrics\textquotedblright{}.) Say it once more.
\end{tcolorbox}
\section[\texorpdfstring{\zh{儿歌} (érgē)}{儿歌 (érgē)}]{\zh{儿歌} (érgē) — nursery rhyme; children's song}

\label{lesson:ZH-C504-erge}



\begin{quote}
Before the new one: say the Mandarin for a new year print, then the Mandarin for the words of a song.
\end{quote}

\subsection*{You'll want to know: \zh{儿歌}}
\textbf{\zh{儿歌}} — \emph{érgē} — \textquotedblleft{}nursery rhyme; children's song\textquotedblright{}.

\begin{quote}
\textbf{\zh{这是我小时候最喜欢的儿歌}} \emph{zhè shì wǒ xiǎoshíhou zuì xǐhuan de érgē} — This was my favourite nursery rhyme when I was little.
\end{quote}

The new character, in pieces: \zh{歌} — fourteen strokes in two pieces: on the left \zh{可} \emph{kě} twice, one over the other, ten strokes; on the right four strokes. The next lesson puts it in your hand.

\begin{cousinweb}
\textbf{\zh{儿歌}} (\emph{érgē}) is built from \zh{儿} 'a child' and \zh{歌} 'a song', both inherited from Middle Chinese, \emph{nye} and \emph{ka} (Guangyun, in Baxter's transcription), which glosses the first 'an infant' and the second 'the voice of a person singing'.
\end{cousinweb}

\begin{grammarlens}[title={What you've built}]
One more word for saying how it was.
\end{grammarlens}

\subsection*{Your turn}
\begin{itemize}
\raggedright
  \item \emph{Say it:} \emph{érgē}
  \item \emph{Say it:} \emph{érgē}, once more
  \item \emph{Say it:} \emph{zhè shì wǒ xiǎoshíhou zuì xǐhuan de érgē} — This was my favourite nursery rhyme when I was little.
  \item \emph{Recall:} say the Mandarin for a new year print, then the Mandarin for the words of a song, then say \emph{érgē} again
\end{itemize}

\begin{tcolorbox}[breakable,colback=teal!4,colframe=teal!35!black,title={Before you move on}]
What does \zh{儿歌} mean? (\textquotedblleft{}Nursery rhyme; children's song\textquotedblright{}.) Say it once more.
\end{tcolorbox}
\section[\texorpdfstring{\zh{歌} (gē)}{歌 (gē)}]{\zh{歌} — kě twice, then four strokes}

\label{lesson:ZH-W504-ge-song}



\begin{quote}
The word first, then the character inside it.

\begin{itemize}
\raggedright
  \item \emph{Recall:} say \emph{érgē} — \textbf{R1}, one lesson back, before it decays
\end{itemize}

\begin{quote}
\zh{可} \zh{儿}
\end{quote}

\emph{Read it:} each one aloud, with nothing said first — \emph{kě}, \emph{ér}
\end{quote}

\subsection*{Script — kě twice, then four strokes}
\hlblockfigure{\detokenize{figures/ZH-W504-ge-song-filmstrip.pdf}}{How \zh{歌} is written, stroke by stroke}

\begin{quote}
\zh{歌}
\end{quote}

\textbf{Two pieces, left before right.}

\textbf{The left piece, ten strokes: \zh{可} \emph{kě} twice, one over the other.} The upper \zh{可} is small: a bar, \zh{口}, and a vertical with no hook. The lower \zh{可} is wider: a long bar, \zh{口}, and a vertical that hooks to the left.

\textbf{The right piece, four strokes}: a short stroke falling to the left, a bar that turns down to the left, a long stroke falling to the left, and a long stroke falling to the right.

5 + 5 + 4 = fourteen.

\subsection*{Your turn}
\begin{itemize}
\raggedright
  \item \emph{Write it:} \zh{歌}, all fourteen, with the model in front of you
  \item \emph{Cover:} write it again from memory, and count the pieces, not the strokes
\end{itemize}

\begin{tcolorbox}[breakable,colback=teal!4,colframe=teal!35!black,title={Before you move on}]
How many strokes does \zh{歌} have? (\textbf{Fourteen}.) What is the left piece? (\textbf{\zh{可}, twice}.) Which of the two has a hook? (\textbf{The lower one}.)

Source: \href{https://raw.githubusercontent.com/chanind/hanzi-writer-data/68d10a4b21150cae5e1ebbd223eed289cf32d90c/data/\%E6\%AD\%8C.json}{Hanzi Writer Data \zh{歌}, pinned stroke paths and medians}.
\end{tcolorbox}
\section[\texorpdfstring{\zh{国歌} (guógē)}{国歌 (guógē)}]{\zh{国歌} (guógē) — national anthem}

\label{lesson:ZH-C504-guoge}



\begin{quote}
Before the new one: say the Mandarin for the words of a song, then the Mandarin for a nursery rhyme.
\end{quote}

\subsection*{You'll want to know: \zh{国歌}}
\textbf{\zh{国歌}} — \emph{guógē} — \textquotedblleft{}national anthem\textquotedblright{}.

\begin{quote}
\emph{tīng guógē de shíhou wǒ gǎnjué hěn jīdòng} — When I hear the national anthem, I feel very moved.
\end{quote}

\begin{cousinweb}
\textbf{\zh{国歌}} (\emph{guógē}) is built from \zh{国} 'a country' and \zh{歌} 'a song', both inherited from Middle Chinese, \emph{kwok} and \emph{ka} (Guangyun, in Baxter's transcription).
\end{cousinweb}

\begin{grammarlens}[title={What you've built}]
One more word for saying how it was.
\end{grammarlens}

\subsection*{Your turn}
\begin{itemize}
\raggedright
  \item \emph{Say it:} \emph{guógē}
  \item \emph{Say it:} \emph{guógē}, once more
  \item \emph{Say it:} \emph{tīng guógē de shíhou wǒ gǎnjué hěn jīdòng} — When I hear the national anthem, I feel very moved.
  \item \emph{Recall:} say the Mandarin for the words of a song, then the Mandarin for a nursery rhyme, then say \emph{guógē} again
\end{itemize}

\begin{tcolorbox}[breakable,colback=teal!4,colframe=teal!35!black,title={Before you move on}]
What does \zh{国歌} mean? (\textquotedblleft{}National anthem\textquotedblright{}.) Say it once more.
\end{tcolorbox}
\section[míngē]{míngē — folk song}

\label{lesson:ZH-C504-minge}



\begin{quote}
Before the new one: say the Mandarin for a nursery rhyme, then the Mandarin for the national anthem.
\end{quote}

\subsection*{You'll want to know: míngē}
\emph{míngē} — \textquotedblleft{}folk song\textquotedblright{}.

\begin{quote}
\emph{wǒ zuì xǐhuan tīng Zhōngguó míngē} — Chinese folk songs are my favourite.
\end{quote}

Its characters come later; for now it is a word for the ear.

\begin{cousinweb}
\emph{míngē} is built from \emph{mín} 'the people' and \zh{歌} 'a song', both inherited from Middle Chinese, \emph{mjin} and \emph{ka} (Guangyun, in Baxter's transcription), which glosses the first 'the common people'.
\end{cousinweb}

\begin{grammarlens}[title={What you've built}]
One more word for saying how it was.
\end{grammarlens}

\subsection*{Your turn}
\begin{itemize}
\raggedright
  \item \emph{Say it:} \emph{míngē}
  \item \emph{Say it:} \emph{míngē}, once more
  \item \emph{Say it:} \emph{wǒ zuì xǐhuan tīng Zhōngguó míngē} — Chinese folk songs are my favourite.
  \item \emph{Recall:} say the Mandarin for a nursery rhyme, then the Mandarin for the national anthem, then say \emph{míngē} again
\end{itemize}

\begin{tcolorbox}[breakable,colback=teal!4,colframe=teal!35!black,title={Before you move on}]
What does míngē mean? (\textquotedblleft{}Folk song\textquotedblright{}.) Say it once more.
\end{tcolorbox}
\section[yǎnchàng]{yǎnchàng — to perform a song (on stage)}

\label{lesson:ZH-C504-yanchang-sing}



\begin{quote}
Before the new one: say the Mandarin for the national anthem, then the Mandarin for a folk song.
\end{quote}

\subsection*{You'll want to know: yǎnchàng}
\emph{yǎnchàng} — \textquotedblleft{}to perform a song (on stage)\textquotedblright{}.

\begin{quote}
\emph{nàge gēshǒu lái wǒmen xuéxiào yǎnchàng} — That singer is coming to sing at our school.
\end{quote}

Its characters come later; for now it is a word for the ear.

\begin{cousinweb}
\emph{yǎnchàng} is built from \emph{yǎn} 'to perform' and \emph{chàng} 'to sing', both inherited from Middle Chinese, \emph{yenX} and \emph{tsyhangH} (Guangyun, in Baxter's transcription), which glosses the first 'to spread wide' and the second 'to strike up a song'; the sense 'to perform' is not given for the first there.
\end{cousinweb}

\begin{grammarlens}[title={What you've built}]
One more word for saying how it was.
\end{grammarlens}

\subsection*{Your turn}
\begin{itemize}
\raggedright
  \item \emph{Say it:} \emph{yǎnchàng}
  \item \emph{Say it:} \emph{yǎnchàng}, once more
  \item \emph{Say it:} \emph{nàge gēshǒu lái wǒmen xuéxiào yǎnchàng} — That singer is coming to sing at our school.
  \item \emph{Recall:} say the Mandarin for the national anthem, then the Mandarin for a folk song, then say \emph{yǎnchàng} again
\end{itemize}

\begin{tcolorbox}[breakable,colback=teal!4,colframe=teal!35!black,title={Before you move on}]
What does yǎnchàng mean? (\textquotedblleft{}To perform a song (on stage)\textquotedblright{}.) Say it once more.
\end{tcolorbox}
